\subsection{Модуль lab5}

\subsubsection{EulerianCycle: проверка, достройка и построение цикла}

\textbf{Вход:} \texttt{Graph<int> const\&}~-- неориентированный граф.

\textbf{Выход:} модифицированный граф после \texttt{makeEulerian};
\texttt{optional<vector<int>>}~-- последовательность вершин эйлерова
цикла (или путь, если граф полуэйлеров).

\textbf{Описание:} Вспомогательные методы: \texttt{degree(v)} суммирует
строку матрицы смежности; \texttt{oddVertices()} возвращает список
вершин с нечётной степенью; \texttt{isConnected()} делает BFS и проверяет
достижимость всех ненулевых вершин; \texttt{isBridge(u, v)} временно
убирает ребро и проверяет через BFS, осталась ли связность.

Граф является эйлеровым, если он связен и все степени чётны:

\begin{lstlisting}[style=cstyle, caption={Проверка эйлеровости (euler.cpp)}]
bool EulerianCycle::isEulerian() const {
    return isConnected() && oddVertices().empty();
}
\end{lstlisting}

Метод \texttt{makeEulerian()} итеративно устраняет нечётные степени.
Для каждой пары нечётных вершин $(u, v)$ выполняется одно из трёх
действий: добавить ребро, если его нет; удалить ребро, если оно не мост;
пропустить пару, если ребро является мостом.

\begin{lstlisting}[style=cstyle, caption={Достройка до эйлерового графа (euler.cpp)}]
for (size_t i = 0; i < odd.size() && !fixed; i++) {
    for (size_t j = i + 1; j < odd.size() && !fixed; j++) {
        int u = odd[i], v = odd[j];
        if (!g.hasEdge(u, v)) {
            g.addEdge(u, v, 1); g.addEdge(v, u, 1);
            addedEdges.emplace_back(u, v);
            fixed = true;
        } else if (!isBridge(u, v)) {
            g.removeEdge(u, v); g.removeEdge(v, u);
            fixed = true;
        }
    }
}
\end{lstlisting}

Метод \texttt{buildCycle()} строит эйлеров цикл алгоритмом Хирхольцера
на основе стека. Пройденные рёбра обнуляются в рабочей копии матрицы.

\begin{lstlisting}[style=cstyle, caption={Построение эйлерова цикла (euler.cpp)}]
S.push(start);
while (!S.empty()) {
    int v = S.top(), u = -1;
    for (int j = 0; j < n; j++)
        if (adj[v][j]) { u = j; break; }
    if (u == -1) { S.pop(); circuit.push_back(v); }
    else { S.push(u); adj[v][u] = adj[u][v] = 0; }
}
reverse(circuit.begin(), circuit.end());
\end{lstlisting}

\textbf{Сложность:} \texttt{makeEulerian}~-- $O(V^2 \cdot (V + E))$
(проверка моста через BFS на каждую пару нечётных вершин);
\texttt{buildCycle}~-- $O(V + E)$ по времени; $O(V + E)$ по памяти.

\subsubsection{CycleBasis: фундаментальная система циклов}

\textbf{Вход:} \texttt{vector<Edge> const\&}~-- рёбра МОД (из Краскала),
\texttt{Graph<int> const\&}~-- исходный граф.

\textbf{Выход:} \texttt{vector<EdgeSet>}~-- система фундаментальных циклов;
\texttt{EdgeSet}~-- результат симметрической разности выбранных циклов.

\textbf{Описание:} Метод \texttt{compute()} перебирает все рёбра исходного
графа. Для каждого ребра $(i,j)$, не входящего в МОД (хорды), вызывается
\texttt{pathInMST(i, j)}~-- BFS по остову, возвращающий путь.
Из этого пути формируется множество рёбер; к нему добавляется хорда~--
получается фундаментальный цикл $Z_e$.

\begin{lstlisting}[style=cstyle, caption={Построение фундаментального цикла для хорды (cycle\_basis.cpp)}]
auto path = pathInMST(i, j, mst);
EdgeSet cycle;
for (int k = 0; k + 1 < (int)path.size(); k++) {
    int u = min(path[k], path[k+1]);
    int v = max(path[k], path[k+1]);
    cycle.insert({u, v});
}
cycle.insert({i, j});
basis.push_back(cycle);
\end{lstlisting}

Метод \texttt{interactiveSymDiff()} предлагает пользователю выбрать
номера циклов и последовательно вычисляет их симметрическую разность:

\begin{lstlisting}[style=cstyle, caption={Симметрическая разность циклов (cycle\_basis.cpp)}]
EdgeSet result = basis[indices[0]];
for (size_t i = 1; i < indices.size(); i++)
    result = symDiff(result, basis[indices[i]]);
\end{lstlisting}

\texttt{symDiff(a, b)} возвращает рёбра, входящие ровно в одно из двух
множеств.

\textbf{Сложность:} построение базиса~-- $O(E \cdot (V + E))$
(BFS по остову на каждую хорду);
симметрическая разность~-- $O(E \log E)$ на каждую операцию;
$O(E)$ по памяти.
